%% ---------------------------------------------------------------------
\section{Espacios normados: medir tamaños}
\label{sec:pm-normas}
%% ---------------------------------------------------------------------

Una métrica responde ``¿qué tan lejos están $x$ y $y$?''. En un espacio vectorial ---donde se puede sumar y multiplicar por escalares--- tiene sentido una pregunta más básica: ``¿qué tan grande es $x$?''.

\begin{definicion}[label=def:pm-norma]{Norma y espacio normado}
Sea $X$ un espacio vectorial real. Una función $\norm{\cdot}:X\to\R$ es una \clave{norma} si para todo $x,y\in X$ y todo $\alpha\in\R$:
\begin{enumerate}[label=(N\arabic*),leftmargin=3em]
  \item $\norm x\ge0$, y $\norm x=0$ si y solo si $x=0$ \hfill(definida positiva)
  \item $\norm{\alpha x}=\abs\alpha\,\norm x$ \hfill(homogeneidad)
  \item $\norm{x+y}\le\norm x+\norm y$ \hfill(desigualdad triangular)
\end{enumerate}
El par $(X,\norm{\cdot})$ es un \clave{espacio normado}.
\end{definicion}

Toda norma induce una métrica, $d(x,y)\defeq\norm{x-y}$: la distancia entre dos elementos es el tamaño del desplazamiento que lleva de uno al otro. Las propiedades (M1)--(M4) salen de (N1)--(N3); por ejemplo, $\norm{x-y}=\norm{(x-z)+(z-y)}\le\norm{x-z}+\norm{z-y}$. Lo contrario no es cierto: la métrica de Hamming no viene de ninguna norma (las palabras no se suman).

En $\R^n$, las métricas $d_p$ de la sección anterior vienen de las \clave{normas-$p$}
\begin{equation}
  \norm x_p=\Bigl(\sum_{i=1}^n\abs{x_i}^p\Bigr)^{1/p}\quad(1\le p<\infty),\qquad \norm x_\infty=\max_i\abs{x_i} .
  \label{eq:pm-normasp}
\end{equation}
Para $x=[3,4]\T$: $\norm x_1=7$, $\norm x_2=5$ y $\norm x_\infty=4$. Cuando no se indica el subíndice, $\norm x$ es la norma euclidiana $\norm x_2$.

\script{cap03_preliminares/C3_Ej2_Normas.m}
\begin{figure}[ht]
\centering
\includegraphics[width=\linewidth]{cap03/f_bolas}
\caption{Izquierda: la ``esfera unitaria'' $\{x:\norm x_p=1\}$ es un rombo para $p=1$, un círculo para $p=2$ y un cuadrado para $p=\infty$. Para $p=1/2$ la figura no es convexa: no es una norma. Derecha: en $C[0,1]$ la sucesión $t^k$ tiene norma del máximo $1$ y norma $\mathcal{L}_1$ que tiende a cero.}
\label{fig:pm-bolas}
\end{figure}

\begin{trampa}[label=tr:pm-pmenor]{La fórmula no hace la norma}
\plausible{$\norm x_p=\bigl(\sum_i\abs{x_i}^p\bigr)^{1/p}$ es una norma para todo $p>0$; para $p=1/2$ simplemente pesa más las componentes pequeñas.}
Para $p<1$ falla la desigualdad triangular. Con $a=[1,0]\T$, $b=[0,1]\T$ y $p=1/2$: $\norm a=\norm b=1$, pero $\norm{a+b}=(1+1)^2=4>2$. Geométricamente, la ``bola unitaria'' de la \cref{fig:pm-bolas} no es convexa: el segmento entre dos puntos de la bola se sale de ella.
\correcto{$\norm\cdot_p$ es una norma si y solo si $p\ge1$ (la triangular para $p\ge1$ es la desigualdad de Minkowski). La bola unitaria de cualquier norma es convexa y simétrica respecto al origen.}
\end{trampa}

\codigorepo{cap03_preliminares/C3_Ej2_Normas.m}{pmedio}{$p=1/2$ viola la desigualdad triangular}

%% ---------------------------------------------------------------------
\subsection{Equivalencia de normas en \texorpdfstring{$\R^n$}{Rn}}
%% ---------------------------------------------------------------------

Los valores cambian con la norma, pero en $\R^n$ nunca se alejan demasiado entre sí.

\begin{teorema}[label=thm:pm-equivalencia]{Equivalencia de normas en dimensión finita}
Si $\norm\cdot_\alpha$ y $\norm\cdot_\beta$ son dos normas en $\R^n$, existen constantes $c_1,c_2>0$ tales que
\[
  c_1\norm x_\alpha\le\norm x_\beta\le c_2\norm x_\alpha\qquad\text{para todo }x\in\R^n .
\]
En particular,
\begin{equation}
  \norm x_2\le\norm x_1\le\sqrt n\,\norm x_2,\qquad
  \norm x_\infty\le\norm x_2\le\sqrt n\,\norm x_\infty,\qquad
  \norm x_\infty\le\norm x_1\le n\,\norm x_\infty .
  \label{eq:pm-equiv}
\end{equation}
\end{teorema}

\begin{proof}
Las tres desigualdades de \eqref{eq:pm-equiv}. Sea $m=\max_i\abs{x_i}=\norm x_\infty$. (i) $m\le\sum_i\abs{x_i}\le n\,m$. (ii) $m^2\le\sum_ix_i^2\le n\,m^2$; tome raíz. (iii) $\norm x_1^2=\sum_i\abs{x_i}^2+\sum_{i\ne j}\abs{x_i}\abs{x_j}\ge\norm x_2^2$; y, por Cauchy--Schwarz (\cref{thm:pm-cs}) con el vector $[1,\dots,1]\T$, $\norm x_1=\sum_i1\cdot\abs{x_i}\le\sqrt n\,\norm x_2$. Las cotas son ajustadas: $x=[1,0,\dots,0]\T$ da la igualdad a la izquierda y $x=[1,1,\dots,1]\T$ a la derecha.
El caso general: basta compararlas con $\norm\cdot_2$. Toda norma es continua respecto a $\norm\cdot_2$, porque $\abs{\norm x_\alpha-\norm y_\alpha}\le\norm{x-y}_\alpha\le\bigl(\max_i\norm{e_i}_\alpha\bigr)\norm{x-y}_1\le\bigl(\max_i\norm{e_i}_\alpha\bigr)\sqrt n\norm{x-y}_2$. La esfera $\{\norm x_2=1\}$ es compacta (\cref{sec:pm-conjuntos}), así que $\norm\cdot_\alpha$ alcanza en ella un mínimo $c_1>0$ (no es cero porque $x\neq0$) y un máximo $c_2$. Por homogeneidad, $c_1\norm x_2\le\norm x_\alpha\le c_2\norm x_2$ para todo $x$.
\end{proof}

¿Para qué sirve? Para elegir la norma que haga más fácil una cuenta. Si se prueba que $\norm{x(t)}_\infty\to0$, también $\norm{x(t)}_2\to0$; si se prueba que $\norm{x(t)}_1\le ce^{-t}$, entonces $\norm{x(t)}_2\le ce^{-t}$. Las propiedades \emph{cualitativas} ---acotado, convergente, estable--- no dependen de la norma en $\R^n$; las constantes sí. Con $x=[3,4]\T$: $5\le7\le5\sqrt2\approx7.07$.

\begin{cuidado}[Supremo y máximo]
El \emph{supremo} de un conjunto $S\subset\R$ es su menor cota superior; el \emph{máximo} es un elemento de $S$ que es cota superior. $S_1=[0,1]$ tiene $\sup S_1=\max S_1=1$; $S_2=(0,1)$ tiene $\sup S_2=1$ y no tiene máximo. Por eso las normas de señales se escriben con $\sup$: la señal $1-e^{-t}$, $t\ge0$, nunca llega a $1$, y su norma del máximo es $1$.
\end{cuidado}

%% ---------------------------------------------------------------------
\subsection{Normas de señales}
%% ---------------------------------------------------------------------

Una trayectoria $x:[0,\infty)\to\R^n$ es un elemento de un espacio de funciones. Se puede medir el tamaño del estado en un instante, $\norm{x(t)}$, o el tamaño de \emph{toda} la trayectoria:
\begin{equation}
  \norm x_{\mathcal{L}_p}=\Bigl(\int_0^\infty\norm{x(t)}^p\,\dif t\Bigr)^{1/p}\quad(1\le p<\infty),\qquad
  \norm x_{\Linf}=\sup_{t\ge0}\norm{x(t)} .
  \label{eq:pm-Lp}
\end{equation}
\notamargen{Para funciones medibles generales, $\mathcal{L}_p$ se define con la integral de Lebesgue y sus elementos son clases de funciones que coinciden casi en todas partes. En el curso basta pensar en señales continuas a trozos.}
Se dice que $x\in\mathcal{L}_p$ si esa cantidad es \emph{finita}. $\norm x_{\mathcal{L}_1}$ es la magnitud acumulada, $\norm x_{\Ldos}^2$ es una \emph{energía} (si $x$ es una corriente por un resistor de $1\,\Omega$, es la energía disipada) y $\norm x_{\Linf}$ es el pico. Pertenecer a uno de estos espacios no implica pertenecer a otro (\cref{fig:pm-senales}): $1/(1+t)\in\Ldos\cap\Linf$ pero $\notin\mathcal{L}_1$; $\sin t\in\Linf$ solamente; $t^{-1/2}$ en $(0,1)$ (y cero después) está en $\mathcal{L}_1$ pero no en $\Ldos$ ni en $\Linf$. Estas normas vuelven en el capítulo~11, donde la energía de una señal es la que entra por un puerto.

\begin{figure}[ht]
\centering
\includegraphics[width=\linewidth]{cap03/f_senales}
\caption{Cuatro señales y los espacios $\mathcal{L}_p[0,\infty)$ a los que pertenecen. Una señal acotada no tiene por qué tener energía finita, y una con área finita no tiene por qué ser acotada.}
\label{fig:pm-senales}
\end{figure}

\begin{trampa}[label=tr:pm-dimension]{La equivalencia de normas es un privilegio de $\R^n$}
\plausible{Todas las normas son equivalentes; si una señal es pequeña en $\mathcal{L}_1$, es pequeña en la norma del máximo.}
En espacios de funciones es falso. En $C[0,1]$, $f_k(t)=t^k$ cumple $\norm{f_k}_\infty=1$ para todo $k$, mientras que $\norm{f_k}_{\mathcal{L}_1}=1/(k+1)\to0$ (\cref{fig:pm-bolas}, derecha). No existe $c$ con $\norm f_\infty\le c\norm f_{\mathcal{L}_1}$: la sucesión ``converge a cero'' en una norma y no en la otra. La prueba del \cref{thm:pm-equivalencia} usó que la esfera unitaria de $\R^n$ es compacta, y eso falla en dimensión infinita.
\correcto{En $\R^n$ se puede cambiar de norma libremente (con constantes). Con señales hay que decir \emph{en qué norma} es pequeña una señal o converge una sucesión, porque las conclusiones cambian.}
\end{trampa}

\codigorepo{cap03_preliminares/C3_Ej2_Normas.m}{funciones}{$\norm{t^k}_\infty=1$ mientras $\norm{t^k}_{\mathcal{L}_1}\to0$}

%% ---------------------------------------------------------------------
\subsection{Normas de matrices}
%% ---------------------------------------------------------------------

Una matriz $A\in\R^{m\times n}$ transforma vectores: $x\mapsto Ax$. Su tamaño natural es cuánto puede amplificar.

\begin{definicion}[label=def:pm-inducida]{Norma inducida de una matriz}
\[
  \norm A_p=\sup_{x\neq0}\frac{\norm{Ax}_p}{\norm x_p}=\max_{\norm x_p=1}\norm{Ax}_p .
\]
Para $p=1,2,\infty$ hay fórmulas cerradas:
\[
  \norm A_1=\max_j\sum_{i=1}^m\abs{a_{ij}}\ \ \text{(columnas)},\qquad
  \norm A_2=\sqrt{\lmax(A\T A)},\qquad
  \norm A_\infty=\max_i\sum_{j=1}^n\abs{a_{ij}}\ \ \text{(filas)} .
\]
\end{definicion}

\begin{proposicion}[label=prop:pm-matrices]{Propiedades de las normas inducidas}
Para $A\in\R^{m\times n}$ y $B\in\R^{n\times\ell}$:
\begin{enumerate}
  \item $\norm{Ax}_p\le\norm A_p\norm x_p$ \quad(es la definición: ninguna dirección amplifica más);
  \item $\norm{AB}_p\le\norm A_p\norm B_p$ \quad(submultiplicativa);
  \item $\tfrac1{\sqrt n}\norm A_\infty\le\norm A_2\le\sqrt m\,\norm A_\infty$ y $\tfrac1{\sqrt m}\norm A_1\le\norm A_2\le\sqrt n\,\norm A_1$;
  \item $\norm A_2\le\sqrt{\norm A_1\norm A_\infty}$.
\end{enumerate}
\end{proposicion}

\begin{ejemplo}[label=ej:pm-matriz]{Tres normas de la misma matriz}
Calcule $\norm A_1$, $\norm A_\infty$ y $\norm A_2$ para $A=\begin{bmatrix}1&2\\-3&4\end{bmatrix}$, y compare $\norm A_2$ con el mayor $\abs{\lambda}$.
\solucion Columnas: $1+3=4$ y $2+4=6$, luego $\norm A_1=6$. Filas: $1+2=3$ y $3+4=7$, luego $\norm A_\infty=7$. $A\T A=\begin{bmatrix}10&-10\\-10&20\end{bmatrix}$ tiene valores propios $15\pm5\sqrt5$, así que $\norm A_2=\sqrt{15+5\sqrt5}\approx5.117\le\sqrt{6\cdot7}\approx6.48$. Los valores propios de $A$ son $\frac{5\pm j\sqrt{15}}{2}$, con módulo $\sqrt{10}\approx3.16$: el \emph{radio espectral} $\rho(A)=\max\abs{\lambda_i}$ es menor que $\norm A_2$. La \cref{fig:pm-matriz} muestra por qué: el círculo unitario se transforma en una elipse cuyo semieje mayor es $\norm A_2$.
\end{ejemplo}

\codigorepo{cap03_preliminares/C3_Ej3_NormasMatrices.m}{inducidas}{Las tres normas inducidas, con fórmula y con \cmd{norm}}

\begin{figure}[ht]
\centering
\includegraphics[width=\linewidth]{cap03/f_matriz}
\caption{Izquierda: $A$ convierte el círculo unitario en una elipse; el semieje mayor es $\norm A_2$, mayor que $\rho(A)$ (círculo discontinuo). Derecha: la matriz nilpotente $N$ tiene $\rho(N)=0$, pero aplasta el círculo en un segmento de longitud $2$: $\norm N_2=1$.}
\label{fig:pm-matriz}
\end{figure}

\begin{trampa}[label=tr:pm-radio]{La norma 2 no es el mayor valor propio}
\plausible{$\norm A_2=\max_i\abs{\lambda_i(A)}$: la amplificación máxima la da el valor propio más grande.}
Solo es cierto si $A$ es normal ($A\T A=AA\T$), por ejemplo simétrica. En general $\rho(A)\le\norm A$ para \emph{cualquier} norma inducida (si $Av=\lambda v$, entonces $\abs\lambda\norm v=\norm{Av}\le\norm A\norm v$), pero la desigualdad puede ser estricta y enorme. $N=\begin{bmatrix}0&1\\0&0\end{bmatrix}$ tiene $\rho(N)=0$ y $\norm N_2=1$.
\correcto{Use $\norm A_2=\sqrt{\lmax(A\T A)}$, el mayor valor singular. La diferencia importa en estabilidad: una matriz Hurwitz puede amplificar mucho antes de decaer ($\norm{e^{At}}$ crece transitoriamente si $A$ no es normal), y las funciones de Lyapunov cuadráticas del capítulo~7 son la manera correcta de manejarlo.}
\end{trampa}

\codigorepo{cap03_preliminares/C3_Ej3_NormasMatrices.m}{espectral}{Radio espectral y norma 2}
